\section{Linear Algebra}

\subsection{Linear algebra}

A scratch chapter for linear-algebra notes --- enough structure to flesh out later, and
enough variety (numbered equations, a table, admonitions, and interleaved Lean) to show
the book template in action. The runnable code lives in
\href{/LinearAlgebra-21b7e76089afe77d1ab624eef2fc49b2.lean}{\texttt{LearningLean/LinearAlgebra.lean}}.

\begin{framed}
\textbf{Note}\\
This chapter is a work in progress. Sections marked \textbf{To flesh out} are placeholders.
\end{framed}

\subsubsection{Vectors}

A vector in $\R^n$ is an ordered list of $n$ real numbers. We write column vectors as
$\vec{u} = (u_1, \dots, u_n)$ with each $u_i \in \R$.

The \textbf{dot product} of $\vec{u}, \vec{v} \in \R^n$ is the scalar

\begin{equation}
\label{eq-dot}
\vec{u} \cdot \vec{v} = \sum_{i=1}^{n} u_i\, v_i .
\end{equation}

For a concrete computation, take $\vec{u} = (1, 2, 3)$ and $\vec{v} = (4, 5, 6)$. Then by
equation (\ref{eq-dot}),

\begin{equation}
\vec{u} \cdot \vec{v} = 1\cdot 4 + 2\cdot 5 + 3\cdot 6 = 32 .
\end{equation}

In Lean, a vector is a function \texttt{Fin n \rightarrow ℝ}, and \texttt{![\dots]} is notation for one. The dot product and
the matrix operations below follow the spellings used in Mathlib \citep{mathlib2020}:\footnote{Mathlib writes the dot product as \texttt{⬝ᵥ} (\texttt{{\textbackslash}cdotv}) rather than reusing \texttt{·}, to keep
it distinct from scalar multiplication.}

\begin{verbatim}
def u : Fin 3 \rightarrow ℝ := ![1, 2, 3]
def v : Fin 3 \rightarrow ℝ := ![4, 5, 6]

#check u ⬝ᵥ v      - - `⬝ᵥ` is Mathlib's dot product, of type ℝ
\end{verbatim}

\subsubsection{Matrices}

An $m \times n$ matrix $A \in \R^{m \times n}$ is a rectangular array of scalars. For example,

\begin{equation}
A = \begin{bmatrix} 1 & 2 \\ 3 & 4 \end{bmatrix} \in \R^{2 \times 2}.
\end{equation}

The \textbf{matrix--vector product} $A\vec{x}$ takes a vector $\vec{x} \in \R^n$ to a vector in
$\R^m$. With $\vec{x} = (5, 6)$:

\begin{equation}
A\vec{x}
= \begin{bmatrix} 1 & 2 \\ 3 & 4 \end{bmatrix}\begin{bmatrix} 5 \\ 6 \end{bmatrix}
= \begin{bmatrix} 1\cdot 5 + 2\cdot 6 \\ 3\cdot 5 + 4\cdot 6 \end{bmatrix}
= \begin{bmatrix} 17 \\ 39 \end{bmatrix}.
\end{equation}

A few operations and their Mathlib spellings:

\bigskip\noindent
\begin{tabular}{p{\dimexpr 0.333\linewidth-2\tabcolsep}p{\dimexpr 0.333\linewidth-2\tabcolsep}p{\dimexpr 0.333\linewidth-2\tabcolsep}}
\toprule
Operation & Math & Lean / Mathlib \\
\hline
Dot product & $\vec{u} \cdot \vec{v}$ & \texttt{u ⬝ᵥ v} \\
Matrix--vector product & $A\vec{x}$ & \texttt{A *ᵥ x} \\
Matrix product & $AB$ & \texttt{A * B} \\
Transpose & $A^\top$ & \texttt{Aᵀ} \\
\bottomrule
\end{tabular}

\bigskip

\begin{verbatim}
def A : Matrix (Fin 2) (Fin 2) ℝ := !![1, 2; 3, 4]

#check A *ᵥ ![5, 6]   - - matrix -vector product, of type `Fin 2 \rightarrow ℝ`
\end{verbatim}

\subsubsection{Linear maps}

A \textbf{linear map} $T : \R^n \to \R^m$ is a function satisfying
$T(\alpha\vec{u} + \vec{v}) = \alpha\,T(\vec{u}) + T(\vec{v})$ for all scalars $\alpha$ and
vectors $\vec{u}, \vec{v}$. Every such map is represented by a matrix, and matrix
multiplication corresponds to composition of the maps.

\begin{framed}
\textbf{Note}\\
\textbf{To flesh out:} state and prove that \texttt{Matrix.mulVec A} is a \texttt{LinearMap}, and connect it to
\texttt{Matrix.toLin}.
\end{framed}

\subsubsection{A worked computation}

\textbf{To flesh out.} A step-by-step example --- e.g. solving $A\vec{x} = \vec{b}$ by row
reduction --- narrated alongside the corresponding Lean definitions, so the computation and its
formalization sit side by side.

\subsubsection{Exercises}

\begin{framed}
\textbf{Note}\\
\textbf{To flesh out.} Add exercises here, e.g.:

\begin{enumerate}
\item Show the dot product is symmetric: $\vec{u} \cdot \vec{v} = \vec{v} \cdot \vec{u}$.
\item Prove $(A^\top)^\top = A$ in Lean using \texttt{Matrix.transpose\_transpose}.
\end{enumerate}
\end{framed}